\documentclass[10pt,a4paper]{article}
\usepackage[utf8]{inputenc}
\usepackage[T1]{fontenc}
\usepackage[portuguese]{babel}
\usepackage{geometry}
\geometry{top=1.2cm,bottom=1.2cm,left=1.2cm,right=1.2cm}
\usepackage{multicol}
\usepackage{listings}
\usepackage{xcolor}
\usepackage{array}
\usepackage{booktabs}

% Impede que o LaTeX estique os espacos verticais
\raggedbottom

% Configuracoes do listings para Python
\lstset{
    language=Python,
    basicstyle=\ttfamily\scriptsize,
    keywordstyle=\color{blue!70!black}\bfseries,
    commentstyle=\color{gray!90}\itshape,
    stringstyle=\color{teal},
    showstringspaces=false,
    breaklines=true,
    frame=single,
    rulecolor=\color{gray!30},
    backgroundcolor=\color{gray!5},
    xleftmargin=2pt,
    xrightmargin=2pt,
    aboveskip=4pt,
    belowskip=4pt
}

\begin{document}

\begin{center}
    {\Large \textbf{Cola de Python: Referência Rápida para Lógica}} \\
    \vspace{0.1cm}
    {\small Redes de Computadores -- IFCE Campus Tauá}
\end{center}
\vspace{0.4cm}

\begin{multicols}{2}

\noindent
\begin{minipage}{\linewidth}
\section*{1. Entrada, Saída e Conversão}
Para interagir com o usuário, ler do teclado e exibir dados formatados.
\begin{lstlisting}
# Entrada de dados (lidos como texto/string)
texto = input("Digite algo: ")
# Conversao de tipos
inteiro = int(input("Digite um inteiro: "))
decimal = float(input("Digite um decimal: "))

# Saida com f-strings (chaves {} para variaveis)
print(f"Texto: {texto} | Inteiro: {inteiro}")
# Formatar com casas decimais (Ex: 2 casas)
print(f"Decimal formatado: {decimal:.2f}")
\end{lstlisting}
\end{minipage}
\vspace{0.4cm}

\noindent
\begin{minipage}{\linewidth}
\section*{2. Operações Matemáticas}
Operadores aritméticos básicos para cálculos no Python.
\begin{lstlisting}
a = 10
b = 3

soma = a + b        # 13
sub = a - b         # 7
mult = a * b        # 30
div = a / b         # 3.3333... (divisao real)
div_inteira = a // b # 3 (descarta a parte decimal)
resto = a % b       # 1 (resto da divisao, util p/ par/impar)
potencia = a ** b   # 1000 (10 elevado a 3)
\end{lstlisting}
\end{minipage}
\vspace{0.4cm}

\noindent
\begin{minipage}{\linewidth}
\section*{3. Manipulação de Strings}
Métodos úteis para tratar textos e strings.
\begin{lstlisting}
texto = "  Algoritmos em Python!  "

maiusculo = texto.upper()     # "  ALGORITMOS EM PYTHON!  "
minusculo = texto.lower()     # "  algoritmos em python!  "
limpo = texto.strip()         # "Algoritmos em Python!" (remove espacos)
trocado = texto.replace("Python", "C") # "  Algoritmos em C!  "

# Dividir texto em lista por um separador
frase = "arroz,feijao,carne"
itens = frase.split(",")       # ['arroz', 'feijao', 'carne']

# Juntar lista em texto
junto = "-".join(itens)        # "arroz-feijao-carne"
\end{lstlisting}
\end{minipage}
\vspace{0.4cm}

\noindent
\begin{minipage}{\linewidth}
\section*{4. Decisões e Condições (\texttt{if})}
Executa blocos de código com base em testes lógicos.
\begin{lstlisting}
# Comparadores: ==, !=, >, <, >=, <=
# Operadores logicos: and, or, not

# --- IF SIMPLES (Verificacao unica) ---
nota = 8.5
if nota >= 7.0:
    print("Aprovado!")

# --- IF COMPOSTO (Multiplos caminhos) ---
idade = 17
tem_autorizacao = True

if idade >= 18:
    print("Acesso liberado (Maior de idade)")
elif idade >= 16 and tem_autorizacao:
    print("Acesso liberado (Com autorizacao)")
else:
    print("Acesso negado")
\end{lstlisting}
\end{minipage}
\vspace{0.4cm}

\noindent
\begin{minipage}{\linewidth}
\section*{5. Repetições (Loops: \texttt{for} e \texttt{while})}
Repete um bloco de código de forma controlada.
\begin{lstlisting}
# --- LOOP FOR COM RANGE ---
# range(fim) -> vai de 0 ate fim-1
for i in range(5):  # 0, 1, 2, 3, 4
    print(i)

# range(inicio, fim) -> de inicio ate fim-1
for i in range(2, 6):  # 2, 3, 4, 5
    print(i)

# range(inicio, fim, passo) -> com saltos
for i in range(10, 21, 2):  # 10, 12, 14, 16, 18, 20
    print(i)

# --- LOOP WHILE (Enquanto) ---
contador = 1
while contador <= 5:
    print(contador)
    contador += 1  # Evita loop infinito!
\end{lstlisting}
\end{minipage}
\vspace{0.4cm}

\noindent
\begin{minipage}{\linewidth}
\section*{6. Listas (Vetores de 1 Dimensão)}
Armazena múltiplos valores ordenados.
\begin{lstlisting}
numeros = [10, 20, 30, 40]
primeiro = numeros[0]    # 10
ultimo = numeros[-1]     # 40 (fim da lista)

# Inserir e Modificar
numeros[1] = 25       # [10, 25, 30, 40]
numeros.append(50)    # Final: [10, 25, 30, 40, 50]
numeros.insert(1, 15) # No indice 1: [10, 15, 25, 30, 40, 50]

# Remover
numeros.pop()         # Remove ultimo: [10, 15, 25, 30, 40]
numeros.pop(1)        # Remove do indice 1: [10, 25, 30, 40]
numeros.remove(30)    # Remove valor 30: [10, 25, 40]

# Ordenacao e Inversao
valores = [5, 2, 9, 1]
valores.sort()        # Altera p/ crescente: [1, 2, 5, 9]
valores.sort(reverse=True) # Decrescente: [9, 5, 2, 1]
valores.reverse()     # Inverte a ordem da lista

# Iterar com indices
for i, num in enumerate(valores):
    print(f"Indice {i} tem o valor {num}")
\end{lstlisting}
\end{minipage}
\vspace{0.4cm}

\noindent
\begin{minipage}{\linewidth}
\section*{7. Matrizes (Listas de Listas)}
Coleções organizadas em linhas e colunas (2 dimensões).
\begin{lstlisting}
matriz = [
    [1, 2, 3],  # Linha 0
    [4, 5, 6],  # Linha 1
    [7, 8, 9]   # Linha 2
]

# Acesso: matriz[linha][coluna]
valor = matriz[1][2]  # Linha 1, Coluna 2 -> valor 6
matriz[0][0] = 99     # Altera valor na Linha 0, Coluna 0

# Percorrer linha por linha
for linha in matriz:
    for elemento in linha:
        print(elemento, end="\t")
    print()  # Quebra de linha ao fim de cada linha

# Percorrer usando indices
linhas = len(matriz)
colunas = len(matriz[0])
for i in range(linhas):
    for j in range(colunas):
        print(f"matriz[{i}][{j}] = {matriz[i][j]}")
\end{lstlisting}
\end{minipage}
\vspace{0.4cm}

\noindent
\begin{minipage}{\linewidth}
\section*{8. Dicionários (Chave-Valor)}
Armazena dados associados a rótulos (chaves).
\begin{lstlisting}
aluno = {
    "nome": "Joao Pedro",
    "idade": 20,
    "curso": "ADS"
}
print(aluno["nome"])  # "Joao Pedro"

# Adicionar/Modificar
aluno["media"] = 8.5  # Cria nova chave
aluno["idade"] = 21   # Modifica chave existente

# Remocao
del aluno["curso"]
valor = aluno.pop("idade")

# Percorrer Chaves e Valores
for chave, valor in aluno.items():
    print(f"{chave}: {valor}")
\end{lstlisting}
\end{minipage}
\vspace{0.4cm}

\noindent
\begin{minipage}{\linewidth}
\section*{9. Funções Embutidas Úteis}
Funções globais que facilitam cálculos comuns.

\noindent\begin{tabular}{@{}llp{3cm}@{}}
\toprule
\textbf{Função} & \textbf{Descrição} & \textbf{Exemplo} \\ \midrule
\texttt{len(c)} & Tamanho da coleção & \texttt{len([4, 8, 2])} $\rightarrow$ \texttt{3} \\
\texttt{sum(c)} & Soma elementos & \texttt{sum([4, 8, 2])} $\rightarrow$ \texttt{14} \\
\texttt{min(c)} & Menor elemento & \texttt{min([4, 8, 2])} $\rightarrow$ \texttt{2} \\
\texttt{max(c)} & Maior elemento & \texttt{max([4, 8, 2])} $\rightarrow$ \texttt{8} \\
\texttt{sorted(c)} & Cria lista ordenada & \texttt{sorted([4, 8, 2])} $\rightarrow$ \texttt{[2, 4, 8]} \\
\texttt{abs(n)} & Valor absoluto & \texttt{abs(-15)} $\rightarrow$ \texttt{15} \\
\texttt{round(n,d)}& Arredonda p/ $d$ casas & \texttt{round(3.1415, 2)} $\rightarrow$ \texttt{3.14} \\ \bottomrule
\end{tabular}
\end{minipage}
\vspace{0.4cm}

\noindent
\begin{minipage}{\linewidth}
\section*{10. Padrões de Algoritmos}
Soluções clássicas aplicadas a listas de dados.
\vspace{0.2cm}
\subsection*{A. Acumulador (Soma acumulada)}
Soma valores com base em critérios.
\begin{lstlisting}
valores = [12, 5, 8, 21, 30, 4]
soma_pares = 0
for v in valores:
    if v % 2 == 0:
        soma_pares += v
\end{lstlisting}
\end{minipage}
\vspace{0.3cm}

\noindent
\begin{minipage}{\linewidth}
\subsection*{B. Contador}
Conta ocorrências de determinado critério.
\begin{lstlisting}
notas = [6.5, 8.0, 5.5, 9.0, 4.0, 7.0]
aprovados = 0
for nota in notas:
    if nota >= 6.0:
        aprovados += 1
\end{lstlisting}
\end{minipage}
\vspace{0.3cm}

\noindent
\begin{minipage}{\linewidth}
\subsection*{C. Maior/Menor Elemento (Manual)}
Algoritmo de extremos sem usar funções embutidas.
\begin{lstlisting}
idades = [25, 42, 18, 50, 31]
maior_idade = idades[0]
for idade in idades:
    if idade > maior_idade:
        maior_idade = idade
\end{lstlisting}
\end{minipage}
\vspace{0.3cm}

\noindent
\begin{minipage}{\linewidth}
\subsection*{D. Filtragem de Listas}
Gera uma nova lista com elementos filtrados.
\begin{lstlisting}
numeros = [1, 2, 3, 4, 5, 6, 7, 8, 9, 10]
impares = []
for n in numeros:
    if n % 2 != 0:
        impares.append(n)
\end{lstlisting}
\end{minipage}
\vspace{0.4cm}

\noindent
\begin{minipage}{\linewidth}
\section*{11. List Comprehension}
Forma curta de criar e filtrar listas em Python.
\begin{lstlisting}
# --- SINTAXE BASICA ---
# [expressao for item in colecao]
quadrados = [x**2 for x in range(5)] # [0, 1, 4, 9, 16]

# --- SINTAXE COM FILTRO (IF) ---
# [expressao for item in colecao if condicao]
pares = [x for x in range(10) if x % 2 == 0] # [0, 2, 4, 6, 8]

# --- EXEMPLOS TEMATICOS (REDES) ---
# 1. Limpar espacos de strings em uma lista de IPs
hosts = [" 192.168.1.1 ", " 10.0.0.1", "172.16.0.254 "]
ips_limpos = [ip.strip() for ip in hosts]
# Resultado: ['192.168.1.1', '10.0.0.1', '172.16.0.254']

# 2. Filtrar apenas IPs locais (comecam com "192.")
ips = ["192.168.1.5", "10.0.0.3", "192.168.2.10", "8.8.8.8"]
ips_locais = [ip for ip in ips if ip.startswith("192.")]
# Resultado: ['192.168.1.5', '192.168.2.10']
\end{lstlisting}
\end{minipage}

\end{multicols}
\end{document}
